% Generated draft: requires mathematical and source review.
\documentclass[11pt]{article}
\usepackage[margin=25mm]{geometry}
\usepackage{fontspec}
\setmainfont{Latin Modern Roman}
\usepackage{amsmath,amssymb,mathtools,unicode-math}
\setmathfont{Latin Modern Math}
\usepackage{enumitem,needspace,xurl,hyperref}
\setlength{\parindent}{0pt}
\setlength{\parskip}{5pt}
\setlength{\emergencystretch}{3em}
\newcommand{\N}{\mathbb{N}}
\newcommand{\Z}{\mathbb{Z}}
\newcommand{\Q}{\mathbb{Q}}
\newcommand{\R}{\mathbb{R}}
\newcommand{\C}{\mathbb{C}}
\newcommand{\F}{\mathbb{F}}
\newcommand{\A}{\mathbb{A}}
\newcommand{\PP}{\mathbb{P}}
\newcommand{\E}{\mathbb{E}}
\newcommand{\eps}{\varepsilon}
\newcommand{\dd}{\,\mathrm{d}}
\begin{document}
\section*{UnsolvedMath ID: 30004061}
\textbf{Source problem code:} catalogue entry
\subsection{Definitions and mathematical statement}
A closed surface or hypersurface is understood here to be compact and without boundary. Write \(\Gamma(t)\) for the evolving geometric object, \(V^\perp\) for its normal velocity, and \(\mathbf H_{\Gamma(t)}\) for its mean-curvature vector. The geometric evolution equation is
\[
V^\perp=\mathbf H_{\Gamma(t)}.
\]
The surface case has intrinsic dimension \(2\); a higher-dimensional hypersurface has intrinsic dimension \(d\geq 3\) and codimension one. The supplied record does not specify the ambient geometry or the codimension of the two-dimensional surface case.

\textbf{Source-stated convergence problem.} Prove convergence of Dziuk's method, or of related evolving finite-element methods, for mean-curvature flow of closed two-dimensional surfaces or higher-dimensional hypersurfaces.

To express the intended limiting assertion schematically, let \(\Gamma_h(t)\) denote a spatial finite-element approximation with mesh parameter \(h\). The requested conclusion has the form
\[
\Gamma_h\longrightarrow\Gamma\qquad\text{as }h\to0,
\]
with a time-step limit also required if a fully discrete method is intended. This display is only a schematic convergence target: the supplied record does not determine a norm, topology, comparison map between discrete and exact surfaces, or time interval. It also does not give the discrete equations, admissible initial approximations, mesh assumptions, element degree, solution regularity, or any coupling between spatial and temporal discretization parameters. Consequently, it does not support a unique, fully quantified convergence theorem. In particular, convergence through singularities or under additional smoothness assumptions must not be inferred from the source wording.

\textbf{Supplied scope qualification.} The record labels the problem partially solved. Its dated assessment states that the historical gap is closed for closed two-dimensional surfaces in several evolving finite-element formulations, and that direct semidiscrete convergence of Dziuk's method is proved for high-order elements, while the broad higher-dimensional and fully general formulation remains incomplete. These are supplied literature assessments, not independently verified conclusions of this draft. The record does not isolate a precise remaining conjecture.
\subsection{Short English statement}
Establish convergence of Dziuk's method or related evolving finite-element approximations to mean-curvature flow for closed two-dimensional surfaces or higher-dimensional hypersurfaces. The source leaves the numerical scheme, hypotheses, convergence topology, and time interval unspecified. Its supplied partially solved status includes results for two-dimensional surfaces and high-order semidiscrete Dziuk approximations; a precise remaining target requires mathematical review.
\subsection{Sources}
\begin{itemize}
\item Supplied local UnsolvedMath registry record, permanent ID \(30004061\), title \emph{Convergence of Finite-Element Mean-Curvature Flow}; collection: Oberwolfach Reports Open Problems; exact source problem code \texttt{catalogue entry}. The record lists proposed year \(2019\) and supplies no named proposer.
\item Source report: \emph{Surface, Bulk, and Geometric Partial Differential Equations: Interfacial, stochastic, non-local and discrete structures} (2020), Oberwolfach Reports, pp.\ 133--207, DOI \texttt{10.4171/owr/2019/3}. Supplied link: \url{https://doi.org/10.4171/owr/2019/3}. The background describes the report as documenting workshop discussions and open problems.
\item Supplied literature reference: Li, convergence of Dziuk's method. \url{https://doi.org/10.1137/20M136935X}.
\item Supplied literature reference: Kovács--Li--Lubich, evolving finite elements. \url{https://doi.org/10.1007/s00211-019-01074-2}.
\item Supplied literature reference: Binz--Kovács, arbitrary codimension. \url{https://doi.org/10.4171/IFB/493}.
\item Literature assessment date: \texttt{2026-08-22}. This is supplied dated curation derived from the OWR record, not independently verified literature assessment. It describes the problem as partially solved, records convergence results for several formulations on closed two-dimensional surfaces and direct semidiscrete convergence of Dziuk's method for high-order elements, and describes the broad higher-dimensional and fully general formulation as incomplete. The source expressly calls for independent verification of time-sensitive status and cited claims.
\item The registry labels its statement \texttt{corrected\_verified} and reports manual semantic, typographic, and LaTeX review. Those are supplied registry qualifications; they do not establish a fully specified mathematical formulation or independent verification of the references.
\end{itemize}
\end{document}
